
%%%%%%%%%%%%%%%%%%%%%%%%%%%%%%%%%%%%%%%%%%%%%%%%%%%%%%%%%%%%%
\section{总体分析}

本题的核心困难在于观测信息的极端稀薄与目标先验的完全缺失。机器狗唯一可用的观测手段是在静止点对单个频道发起一次测向，得到的只是一个叠加了 $\pm1^\circ$ 环境偏差的示向度读数，既没有距离信息，也没有信号强度信息；而目标区域内干扰源的个数在 $10\sim16$ 之间但具体数值未知，每个干扰源占据一个互不相同的频道，频道取自编号 $1\sim20$ 的 20 个信道。这种"只有角度、没有距离"的观测结构决定了全部四个问题的数学骨架：单点观测只能把干扰源约束在一条张角 $2^\circ$ 的楔形之内，必须依靠多点观测的楔形求交才能把无界的方向不确定性收缩为有界的位置不确定性，而交会几何的良窳直接决定定位精度，进而决定机器狗需要走多少冤枉路。

四个问题由此构成一条层层递进的链条。问题一解决纯几何层面的度量问题，即给定若干检测点及其示向度，如何精确计算交会定位区域的直径并判断直径圆能否覆盖该区域；它提供了整个模型的度量工具与精度基准，其结论表明直径圆覆盖率高达 $98.99\%$ 但并非必然成立，失效仅出现在交会角接近 $90^\circ$ 的临界构型上。问题二在只掌握单个检测点示向度的条件下回答第二个检测点应当放在哪里，把几何度量转化为布站决策；它给出的闭式最优基线 $b^{*}=\sqrt2r$ 与数值最优布站为该问题提供了可执行的答案，其核心发现是最优交会角并非 $90^\circ$ 而是 $54.74^\circ$。问题三把问题二的结论嵌入 $N$ 个未知目标的全域搜索任务，在保证全部清除的硬约束下最小化包含移动、频道切换、检测与清除在内的虚拟总时间，本质上是一个带不确定观测的部分可观测序贯决策问题；其核心是证明最少需要 7 个观测点才能保证覆盖，并据此设计"全域粗探—交会定距—迭代逼近—就地清除"的四阶段闭环。问题四在问题三之上引入定向干扰源，其可检测区域由圆盘收缩为半圆盘且朝向未知，使"无信号"这一反馈从确定信息退化为三义信息，需要额外的方位冗余与背向判别机制；其核心结论是零漏测要求观测点在干扰源附近构成"方位包围"，而问题三的七点覆盖网在半径 1200 米以外对定向源必然漏测。

四个问题的技术路线可以概括为三条并行的数学工具线。第一条是凸集几何线，以半平面求交、凸包与旋转卡壳为核心，服务于定位区域的计算与直径的求解；覆盖判据由其延伸出的 Thales 逆定理给出。第二条是估计与优化线，以方位线扰动的误差传播公式为核心，先解析后数值地确定最优布站与候选区域，并进一步导出进入清除精度所需的逼近距离阈值。第三条是组合优化与决策线，以 Kershner 覆盖定理确定观测点数的下界，以 Beardwood--Halton--Hammersley 定理确定行程下界，以方位包围判据确定定向场景下的观测网规模，三条结论共同构成策略评价的参照系，使"策略是否够好"这一问题有了客观的判准而非依赖经验判断。
